\documentclass[12pt]{beamer}

\usepackage[brazil]{babel}
\usepackage[utf8]{inputenc}
\usepackage[T1]{fontenc}
\usepackage{animate}
\usepackage{amsbsy}
\usepackage{amsfonts}
\usepackage{amsmath}
\usepackage{amssymb}
\usepackage{amsthm}
\setbeamertemplate{theorems}[numbered] % to number
\usepackage[toc,page,title,titletoc]{appendix}
\usepackage{dsfont}
\usepackage{esvect}
\usepackage[labelfont=bf]{caption}
%\usepackage{subcaption}
\usepackage{float}
\usepackage[Glenn]{fncychap}%Sonny %Conny %Lenny %Glenn %Renje %Bjarne %Bjornstrup
\usepackage{graphicx}
\usepackage{subfig}
\usepackage{indentfirst}%Para indentar os paragrafos automáticamente
\usepackage{lipsum}
\usepackage{longtable}
\usepackage{mathtools}
\usepackage{listings}%Inserir codigo do R no latex
\usepackage{multirow}
\usepackage{multicol}
\usepackage{csquotes}
\usepackage[maxcitenames=2,terseinits=true,natbib=true, style=authoryear, maxbibnames=99]{biblatex}
\addbibresource{../Referencias/Referencias.bib}
\usepackage[figuresright]{rotating}
\usepackage{spalign}
\usepackage{pgfplots}
\pgfplotsset{compat=1.17}
\usepackage{tikz}
\usepackage{fontawesome}
\usepackage{color, colortbl}
\usepackage{url}
\usepackage{cancel}
\usepackage{accents}
\usepackage{bm}
\usepackage{ragged2e}%para justificar o texto dentro de algum ambiente
\definecolor{Gray}{gray}{0.9}
\definecolor{LightCyan}{rgb}{0.88,1,1}


\usepackage[all]{xy}
\usepackage{hyperref,bookmark}
\hypersetup{
  colorlinks=true,
  linkcolor=blue,
  citecolor=red,
  filecolor=blue,
  urlcolor=blue,
}

\usetheme{Madrid}
\usecolortheme[RGB={193,0,0}]{structure}

%\setbeamertemplate{footline}[frame number]
%\setbeamertemplate{footline}[text line]{%
%  \parbox{\linewidth}{\vspace*{-8pt}\hfill\date{}\hfill\insertshortauthor\hfill\insertpagenumber}}
\beamertemplatenavigationsymbolsempty
\renewcommand{\vec}[1]{\mbox{\boldmath$#1$}}
\newtheorem{Teorema}{Teorema}
\newtheorem{Proposicao}{Proposição}
\newtheorem{definicao}{Definição}
\newtheorem{Corolario}{Corolário}
\newtheorem{Demonstracao}{Demonstração}
\newcommand{\bx}{\ensuremath{\bar{x}}}
\newcommand{\Ho}{\ensuremath{H_{0}}}
\newcommand{\Hi}{\ensuremath{H_{1}}}
\newcommand{\at}[2][]{#1|_{#2}}
\newcommand\xuparrow[1][2ex]{%
   \mathrel{\rotatebox{-90}{$\xleftarrow{\rule{#1}{0pt}}$}}
}
\apptocmd{\frame}{}{\justifying}{} % Allow optional arguments after frame.

\makeatletter
\setbeamertemplate{footline}
{
  \leavevmode%
  \hbox{%
  \begin{beamercolorbox}[wd=.3\paperwidth,ht=2.25ex,dp=1ex,center]{author in head/foot}%
    \usebeamerfont{author in head/foot}\mytext
  \end{beamercolorbox}%
  \begin{beamercolorbox}[wd=.3\paperwidth,ht=2.25ex,dp=1ex,center]{title in head/foot}%
    \usebeamerfont{title in head/foot}\mytextt
  \end{beamercolorbox}%
  \begin{beamercolorbox}[wd=.35\paperwidth,ht=2.25ex,dp=1ex,right]{site in head/foot}%
    \usebeamerfont{site in head/foot}\mytexttt\hspace*{2em}
    \insertframenumber{} / \inserttotalframenumber\hspace*{2ex} 
  \end{beamercolorbox}}%
  \vskip0pt%
}
\makeatother

\providecommand{\arcsin}{} \renewcommand{\arcsin}{\hspace{2pt}\textrm{arcsen}}
\providecommand{\sin}{} \renewcommand{\sin}{\hspace{2pt}\textrm{sen}}
\newcommand{\N}{\rm I\!N}
\newcommand{\I}{\rm I\!I}
\newcommand{\R}{\rm I\!R}
\newcommand{\Sim}{\overset{\text{iid}}{\sim}}
\newcommand{\Lim}{{\displaystyle \lim_{n\to\infty}}}
\newcommand{\LimInf}{{\displaystyle \liminf_{n\to\infty}}}
\newcommand{\rightLim}{\xrightarrow[n\rightarrow\infty]{}}
\newcommand{\Sumi}{{\displaystyle \sum_{i=1}^{n}}}
\newcommand{\Int}{{\displaystyle \int_{-\infty}^{+\infty}}}
\newcommand{\ConvD}{\overset{D}{\rightarrow}}
\newcommand{\ConvP}{\overset{P}{\rightarrow}}
\newcommand{\Prodi}{{\displaystyle \prod_{i=1}^{n}}}
\newcommand{\SetaUP}[2]{\underset{\mathclap{\substack{\xuparrow[30pt] \\ #1}}}{#2}}
%\newcommand{\SetaInclinada}[2]{\underset{\mathclap{\substack{\rotatebox{135}{\xuparrow[30pt] \\ #1}}}}{#2}}
\newcommand{\Home}{\begin{tikzpicture}
\node[scale=2] at (3,4) {\text{Para}~\faHome};
\end{tikzpicture}}
\newcommand{\vecX}{\boldsymbol{X}}
\newcommand{\Implica}[1]{\xRightarrow{#1}}
\newcommand{\SeSe}{\iff}
\newcommand{\EscoreA}{\dfrac{\partial}{\partial\theta}\log{f(x,\theta)}}
\newcommand{\EscoreB}{\dfrac{\partial^{2}}{\partial\theta^{2}}\log{f(x,\theta)}}
\newcommand{\cqd}{\text{cqd}~\blacksquare}
\newcommand{\seqX}{$X_{1},\ldots,X_{n}$}
\newcommand{\seqY}{$Y_{1},\ldots,Y_{n}$}
\newcommand{\tend}[1]{\hbox{\oalign{$\bm{#1}$\crcr\hidewidth$\scriptscriptstyle\bm{\sim}$\hidewidth}}}

%\newtheorem{Teorema}{Teorema}
%\newtheorem{Proposicao}{Proposição}
%\newtheorem{Definicao}{Definição}
%\newtheorem{Corolario}{Corolário}
%\newtheorem{Demonstracao}{Demonstração}

\titlegraphic{\hspace*{8cm}\href{https://fsbmat-ufv.github.io/}{\includegraphics[width=2cm]{figs/mylogo.png}}
}

%Continuar a numeracao em slides diferentes
\newcounter{saveenumi}
\newcommand{\seti}{\setcounter{saveenumi}{\value{enumi}}}
\newcommand{\conti}{\setcounter{enumi}{\value{saveenumi}}}

\resetcounteronoverlays{saveenumi}

% Layout da pagina
\hypersetup{pdfpagelayout=SinglePage}

%Para o \pause funcionar dentro do ambiente align
\makeatletter
\let\save@measuring@true\measuring@true
\def\measuring@true{%
  \save@measuring@true
  \def\beamer@sortzero##1{\beamer@ifnextcharospec{\beamer@sortzeroread{##1}}{}}%
  \def\beamer@sortzeroread##1<##2>{}%
  \def\beamer@finalnospec{}%
}
\makeatother

\title{Inferência Estatística II}
\author{Prof. Fernando de Souza Bastos\texorpdfstring{\\ fernando.bastos@ufv.br}{}}
\institute{Departamento de Estatística\texorpdfstring{\\ Programa de Pós-Graduação em Estatística Aplicada e Biometria}\texorpdfstring{\\ Universidade Federal de Viçosa}{}\texorpdfstring{\\ Campus UFV - Viçosa}{}}
\date{}
\newcommand\mytext{Aula 13}
\newcommand\mytextt{Fernando de Souza Bastos}
\newcommand\mytexttt{\url{https://est711.github.io/}}


\begin{document}
%\SweaveOpts{concordance=TRUE}

\frame{\titlepage}

\begin{frame}{}
\frametitle{\bf Sumário}
\tableofcontents
\end{frame}

\section{Testes Uniformemente Mais Poderosos}
\begin{frame}{}
\begin{block}{}
\justifying
Considere a função densidade de probabilidade (pdf)
\begin{align*}
f(x; \theta) = \begin{cases}
\frac{1}{\theta} e^{-x/\theta}, & 0 < x < \infty \\
0, & \text{Caso Contrário}
\end{cases}
\end{align*}
dos Exercícios 8.1.2 e 8.1.3. Deseja-se testar a hipótese simples $H_0: \theta = 2$ contra a hipótese alternativa composta $H_1: \theta > 2$. Assim, $\Omega = \{\theta: \theta \geq 2\}$. Uma amostra aleatória, $X_1, X_2$, de tamanho $n = 2$ é utilizada e a região crítica é $C = \{(x_1, x_2): 9.5 \leq x_1 + x_2 < \infty\}$. Foi mostrado nos exercícios citados que o nível de significância do teste é aproximadamente $0.05$ e o poder do teste quando $\theta = 4$ é aproximadamente $0.31$. A função poder $\gamma(\theta)$ do teste para todos os $\theta \geq 2$ é:
\end{block}
\end{frame}

\begin{frame}{}
\begin{block}{}
\justifying
\begin{align*}
\gamma(\theta) &= 1 - \int_0^{9.5} \int_0^{9.5-x_2} \frac{1}{\theta^2} e^{-\frac{x_1+x_2}{\theta}} dx_1dx_2 \\
&= \left(\dfrac{\theta+9.5}{\theta}\right)e^{-9.5/\theta}, \quad 2 \leq \theta.
\end{align*}
Por exemplo, $\gamma(2) = 0.05$, $\gamma(4) = 0.31$ e $\gamma(9.5) = \frac{2}{e} \approx 0.74$. É mostrado (Exercício 8.1.3) que o conjunto $C = \{(x_1, x_2): 9.5 \leq x_1 + x_2 < \infty\}$ é uma melhor região crítica de tamanho $0.05$ para testar a hipótese simples $H_0: \theta = 2$ contra cada hipótese simples na hipótese composta $H_1: \theta > 2$.
\end{block}
\end{frame}

\begin{frame}{}
\begin{block}{}
\justifying
%O exemplo anterior ilustra um teste de uma hipótese simples $H_0$ que é o melhor teste de $H_0$ contra todas as hipóteses simples na hipótese alternativa composta $H_1$. Agora definimos uma região crítica, quando ela existe, que é a melhor região crítica para testar uma hipótese simples $H_0$ contra uma hipótese alternativa composta $H_1$. Parece desejável que essa região crítica seja a melhor região crítica para testar $H_0$ contra cada hipótese simples em $H_1$. Ou seja, a função poder do teste que corresponde a essa região crítica deve ser pelo menos tão grande quanto a função poder de qualquer outro teste com o mesmo nível de significância para cada hipótese simples em $H_1$.
\begin{itemize}
	\item \textbf{Resumo do exemplo anterior:}
	\begin{itemize}
		\item exemplo anterior ilustra um teste de uma hipótese simples $H_0$ que é o melhor teste de $H_0$ contra todas as hipóteses simples na hipótese alternativa composta $H_1$.
	\end{itemize}
	
	\item \textbf{Queremos agora definir uma Região Crítica para uma hipótese composta que:}
	\begin{itemize}
		\item Seja a ``melhor'' para testar $H_0$ contra $H_1$.\pause
		\item Melhor região crítica significa:
		\begin{itemize}
			\item A função poder do teste deve ser \textbf{maior ou igual} à de qualquer outro teste para cada hipótese simples em $H_1$.\pause
			\item Deve ter o \textbf{mesmo nível de significância}.
		\end{itemize}
	\end{itemize}
\end{itemize}

\end{block}
\end{frame}

\begin{frame}{}
\begin{definicao}
\justifying
A região crítica $C$ é uma região crítica uniformemente mais poderosa (UMP) de tamanho $\alpha$ para testar a hipótese simples $H_0$ contra uma hipótese alternativa composta $H_1$ se o conjunto $C$ é a melhor região crítica de tamanho $\alpha$ para testar $H_0$ contra cada hipótese simples em $H_1$. Um teste definido por essa região crítica $C$ é chamado de teste uniformemente mais poderoso (UMP), com nível de significância $\alpha$, para testar a hipótese simples $H_0$ contra a hipótese alternativa composta $H_1$.
\end{definicao}
\end{frame}

\begin{frame}{}
\begin{block}{}
\justifying
Como será visto posteriormente, testes uniformemente mais poderosos nem sempre existem. No entanto, quando existem, o teorema de Neyman-Pearson fornece uma técnica para encontrá-los. Alguns exemplos ilustrativos são dados aqui.
\end{block}
\end{frame}

\section{Exemplo 1}
\begin{frame}{Exemplo}
\begin{block}{}
\justifying
Seja $X_1, X_2, \ldots, X_n$ uma amostra aleatória de uma distribuição normal $N(0, \theta)$, onde a variância $\theta$ é um número positivo desconhecido. Será demonstrado que existe um teste uniformemente mais poderoso com nível de significância $\alpha$ para testar a hipótese simples $H_0: \theta = \theta_0$, onde $\theta_0$ é um número positivo fixo, em oposição à hipótese alternativa composta $H_1: \theta > \theta_0$. Assim, $\Omega = \{\theta : \theta \geq \theta_0\}$.

A função de densidade conjunta de $X_1, X_2, \ldots, X_n$ é dada por
\begin{equation}
L(\theta; x_1, x_2, \ldots, x_n) = \frac{1}{(2\pi\theta)^{n/2}} \exp\left(-\frac{1}{2\theta}\sum_{i=1}^{n}x_i^2\right).
\end{equation}
\end{block}
\end{frame}

\begin{frame}{}
\begin{block}{}
\justifying
Seja $\theta_1$ representando um número maior que $\theta_0$, e $k$ denote um número positivo. Seja $C$ o conjunto de pontos onde
\begin{equation}
\frac{L(\theta_0; x_1, x_2, \ldots, x_n)}{L(\theta_1; x_1, x_2, \ldots, x_n)} \leq k,
\end{equation}
isto é, o conjunto de pontos onde
\begin{equation}
\frac{\exp\left(-\frac{1}{2\theta_0}\sum_{i=1}^{n}x_i^2\right)}{\exp\left(-\frac{1}{2\theta_1}\sum_{i=1}^{n}x_i^2\right)} \leq k,
\end{equation}
ou equivalentemente,
\begin{align*}
\sum_{i=1}^{n}x_i^2 \geq \dfrac{2\theta_1\theta_0}{\theta_1-\theta_0}\left(\dfrac{n}{2}\log{\left(\dfrac{\theta_1}{\theta_0}\right)}-\log{k}\right)=c
\end{align*}
\end{block}
\end{frame}

\begin{frame}{}
\begin{block}{}
\justifying
O conjunto $C = \{(x_1, x_2, \ldots, x_n) : \sum_{i=1}^{n}x_i^2 \geq c\}$ é então uma melhor região crítica para testar a hipótese simples $H_0: \theta = \theta_0$ em oposição à hipótese simples $\theta = \theta_1$. Resta determinar $c$ de forma que esta região crítica tenha o tamanho desejado $\alpha$. Se $H_0$ for verdadeira, a variável aleatória
\begin{equation}
\frac{\sum_{i=1}^{n}X_i^2}{\theta_0}
\end{equation}
tem uma distribuição qui-quadrado com $n$ graus de liberdade. Uma vez que $\alpha = P_{\theta_0} \left(\frac{\sum_{i=1}^{n}X_i^2}{\theta_0} \geq \frac{c}{\theta_0}\right)$, $c/\theta_0$ pode ser calculado, por exemplo, usando o código R 

\begin{align*}
    {\texttt qchisq(1 - \alpha, n)}.
\end{align*}

Então, $C = \{(x_1, x_2, \ldots, x_n) : \sum_{i=1}^{n}x_i^2 \geq c\}$ é uma melhor região crítica de tamanho $\alpha$ para testar $H_0: \theta = \theta_0$ em oposição à hipótese $\theta = \theta_1$.


\end{block}
\end{frame}

\begin{frame}{}
\begin{block}{}
\justifying
Além disso, para cada número $\theta_1$ maior que $\theta_0$, o argumento anterior se mantém. Ou seja, $C = \{(x_1, \ldots, x_n) : \sum_{i=1}^{n}x_i^2 \geq c\}$ é uma região crítica uniformemente mais poderosa de tamanho $\alpha$ para testar $H_0: \theta = \theta_0$ em oposição a $H_1: \theta > \theta_0$. Se $x_1, x_2, \ldots, x_n$ denotam os valores experimentais de $X_1, X_2, \ldots, X_n$, então $H_0: \theta = \theta_0$ é rejeitada ao nível de significância $\alpha$, e $H_1: \theta > \theta_0$ é aceita se $\sum_{i=1}^{n}x_i^2 \geq c$; caso contrário, $H_0: \theta = \theta_0$ é aceita.

Se, na discussão anterior, tomarmos $n = 15$, $\alpha = 0,05$ e $\theta_0 = 3$, então as duas hipóteses são $H_0: \theta = 3$ e $H_1: \theta > 3$. Usando R, $c/3$ é calculado por \texttt{qchisq(0.95, 15)} $= 24,996$. Portanto, $c = 74,988$.

\end{block}
\end{frame}

\section{Exemplo 2}
\begin{frame}{Exemplo 2}
\begin{block}{}
\justifying
Seja $X_1, X_2, \ldots, X_n$ uma amostra aleatória de uma distribuição normal $N(\theta, 1)$, onde $\theta$ é desconhecido. Será demonstrado que não existe um teste uniformemente mais poderoso para a hipótese simples $H_0 : \theta = \theta_0$, onde $\theta_0$ é um número fixo, em oposição à hipótese alternativa composta $H_1 : \theta \neq \theta_0$. Assim, $\Omega = \{\theta : -\infty < \theta < \infty\}$.
\end{block}
\pause
\begin{block}{}
	\justifying
	Seja $\theta_1$ um número diferente de $\theta_0$. Seja $k$ um número positivo e considere
	\[
	\dfrac{\frac{1}{(2\pi)^{n/2}} \exp\left(-\frac{1}{2}{\displaystyle \sum_{i=1}^{n}(x_i - \theta_{0})^2}\right) }{\frac{1}{(2\pi)^{n/2}} \exp\left(-\frac{1}{2}{\displaystyle \sum_{i=1}^{n}(x_i - \theta_1)^2}\right)} \leq k.
	\]
\end{block}
\end{frame}

\begin{frame}{}
\begin{block}{}
\justifying
A desigualdade anterior pode ser escrita como
\[
(\theta_{0} - \theta_{1})\sum_{i=1}^{n}x_i \leq \dfrac{2\log k -n(\theta_{1}^{2}-\theta_{0}^{2})}{2}.
\]
\end{block}
\pause
\begin{block}{}
	\justifying
	\begin{itemize}
		\item Se $\theta_{0} > \theta_1$ essa última desigualdade é equivalente a
		\[
		\sum_{i=1}^{n}x_i \leq \dfrac{2\log k -n(\theta_{1}^{2}-\theta_{0}^{2})}{2(\theta_{0} - \theta_{1})},
		\]\pause
		\item Por outro lado, se $\theta_{0} < \theta_1,$ é equivalente a
		\[
		\sum_{i=1}^{n}x_i \geq \dfrac{2\log k -n(\theta_{1}^{2}-\theta_{0}^{2})}{2(\theta_{0} - \theta_{1})}.
		\] 
	\end{itemize}
\end{block}
\end{frame}

\begin{frame}{}
\begin{block}{}
\justifying

%A primeira dessas duas expressões define uma melhor região crítica para testar $H_0 : \theta = \theta_0$ contra a hipótese $\theta = \theta_1$, desde que $\theta_1 > \theta$, enquanto a segunda expressão define uma melhor região crítica para testar $H_0 : \theta = \theta_0$ contra a hipótese $\theta = \theta_1$, desde que $\theta_1 < \theta$. Ou seja, uma melhor região crítica para testar a hipótese simples contra uma hipótese simples alternativa, digamos $\theta = \theta_0 + 1$, não serve como uma melhor região crítica para testar $H_0 : \theta = \theta_0$ contra a hipótese simples alternativa $\theta = \theta_0 - 1$. Por definição, então, não existe um teste uniformemente mais poderoso no caso em consideração.
%
%Deve-se notar que se a hipótese alternativa composta fosse unidirecional, seja $H_1 : \theta > \theta_0$ ou $H_1 : \theta < \theta_0$, um teste uniformemente mais poderoso existiria em cada instância.
\begin{itemize}
	\item \textbf{Definição de Melhor Região Crítica:}
	\begin{itemize}
		\item \textbf{Primeira expressão:} Melhor região crítica para testar $H_0 : \theta = \theta_0$ contra $H_1 : \theta > \theta_0$.\pause
		\item \textbf{Segunda expressão:} Melhor região crítica para testar $H_0 : \theta = \theta_0$ contra $H_1 : \theta < \theta_0$.
	\end{itemize}
	\pause
	\item \textbf{Limitações da Melhor Região Crítica:}
	\begin{itemize}
		\item Uma melhor região crítica para $\theta = \theta_0 + 1$ não serve para testar $H_0 : \theta = \theta_0$ contra $\theta = \theta_0 - 1$.\pause
		\item Portanto, não existe um teste uniformemente mais poderoso neste caso.
	\end{itemize}
	\pause
	\item \textbf{Teste Uniformemente Mais Poderoso:}
	\begin{itemize}
		\item Caso a hipótese alternativa fosse composta unidirecional:
		\begin{itemize}
			\item $H_1 : \theta > \theta_0$ ou $H_1 : \theta < \theta_0$. Um teste uniformemente mais poderoso existiria para cada instância.
		\end{itemize}
	\end{itemize}
\end{itemize}
\end{block}
\end{frame}

\begin{frame}{\Home}
\begin{block}{}
\justifying

\begin{itemize}
    \item \textbf{Exercícios da seção 8.2:} 1,3,5,6,11,13.
\end{itemize}
\end{block}
\nocite{hogg}
\end{frame}

\begin{frame}[allowframebreaks]
\frametitle{\bf Referências}
\printbibliography
\end{frame}


\end{document}

